\chapter{Conclusões}
\label{cap:conclusao}

Este trabalho abordou o problema de otimização do plano de voo de veículos aéreos não tripulados (VANTs) para a coleta de dados de sensores IoT distribuídos em uma área geográfica, com o objetivo de equilibrar a maximização da informação coletada, medida pela \textit{Age of Information} (AoI), e a minimização do consumo energético do veículo. Para isso, propôs-se uma formulação de Programação Linear Inteira Mista (PLIM) com horizonte temporal discretizado em \textit{slots}, na qual a movimentação do VANT é representada sobre um grafo completo de nós (sensores e base) e o custo energético de cada deslocamento, bem como dos períodos de pairar (\textit{hovering}), é estimado por um modelo de potência de asa rotativa.

Os objetivos específicos estabelecidos foram atendidos. O problema foi formulado com tempo discretizado e variáveis de decisão adequadas, contemplando a dinâmica de movimento, a dinâmica de AoI e a restrição de capacidade de bateria. Os objetivos de maximização da AoI coletada e de minimização da energia foram combinados por meio de uma otimização multiobjetivo lexicográfica, evitando a calibração de um peso relativo arbitrário entre as grandezas. A formulação foi implementada no resolvedor Gurobi e avaliada em 36 cenários, variando a quantidade de sensores e o tamanho do mapa, com os parâmetros físicos do DJI Matrice 300 RTK e do Parrot Anafi USA, sendo comparada a heurísticas construtivas de referência (\textit{Nearest Neighbor}, AoI-Greedy e Score-Greedy).

Os experimentos indicaram que o modelo exato é capaz de encontrar soluções de boa qualidade nas escalas avaliadas. Em particular, a análise da variante com revisitas ($K=3$) mostrou que permitir múltiplas coletas por sensor não trouxe melhora estatisticamente significativa na AoI média final (variações inferiores a $\pm 3\%$), ao mesmo tempo em que elevou o consumo energético de forma expressiva em mapas maiores. Conclui-se, portanto, que a restrição de visita única a cada sensor é uma escolha eficiente: ela força o otimizador a diversificar as visitas, obtendo qualidade de AoI equivalente a um custo energético significativamente menor.

Algumas limitações do modelo devem ser destacadas. A discretização adotada conclui cada transição entre nós em um único \textit{slot} a velocidade constante, o que mantém o custo energético linear por aresta, mas implica velocidades nominais irrealistas em mapas de grande escala. Nesses cenários, os valores de energia devem ser interpretados como uma aproximação conservadora do custo de deslocamento. Além disso, o estudo concentrou-se em um único VANT e adotou um modelo de comunicação simplificado, no qual a coleta ocorre apenas quando o veículo paira diretamente sobre o sensor, sem modelar aspectos de camada de enlace. Soma-se a isso o custo computacional crescente do modelo exato: em instâncias maiores, o consumo de tempo e de memória do resolvedor torna-se um obstáculo prático, o que motivou a avaliação de heurísticas construtivas como alternativa de baixo custo (Seção~\ref{sec:escala-heuristicas}).

Como trabalhos futuros, vislumbram-se: (i) a extensão da formulação para múltiplos VANTs cooperativos, explorando ganhos de capacidade e cobertura, (ii) a reformulação do modelo com restrições explícitas de velocidade, desacoplando a velocidade de voo da duração do \textit{slot} e permitindo mais de um \textit{slot} por transição, de modo a representar perfis de voo fisicamente realizáveis e abranger áreas maiores sem comprometer a comparabilidade entre cenários, (iii) o tratamento de instâncias de maior porte por meio de técnicas de decomposição ou matheurísticas, e (iv) a incorporação de modelos de comunicação e de geração de dados mais detalhados, aproximando a formulação de condições operacionais reais.
